\section{Properties of the Wilson flow at small coupling}

The aim in this section partly is to show that the
Wilson flow can be studied straightforwardly in perturbation
theory and partly to check that
the expectation values of local gauge-invariant
observables calculated at positive flow time
are renormalized quantities.

For simplicity the perturbation expansion
is discussed in the continuum theory using
dimensional regularization.
The gauge group is taken to be $\SUn$
and it is assumed that there are $\Nf$ flavours of
massless quarks. As a representative case,
the observable
\begin{equation}
   E=\frac{1}{4}G_{\mu\nu}^aG_{\mu\nu}^a
   \label{eq:Edens}
\end{equation}
is considered and its expectation value is worked out to
next-to-leading order in the gauge coupling.


\subsection{Gauge fixing}

The flow equation~\eqref{eq:floweq} is invariant
under $t$-independent gauge transformations.
$E$ is therefore a gauge-invariant function of the
fundamental field $A_{\mu}$ and its expectation value
can consequently be calculated in any gauge.

In perturbation theory, the gauge invariance of
the flow equation leads to some
technical complications that are better avoided by
considering the modified equation
\begin{equation}
  \dot{B}_{\mu}=D_{\nu}G_{\nu\mu}+\lambda D_{\mu}\partial_{\nu}B_{\nu}.
  \label{eq:modflow}
\end{equation}
For any given value of the gauge parameter $\lambda$,
the solution of eq.~\eqref{eq:modflow} is related to the one at
$\lambda=0$ through
\begin{equation}
  B_{\mu}=\Lambda\!\left.B_{\mu}\right|_{\lambda=0}\Lambda^{-1}
  +\Lambda\partial_{\mu}\Lambda^{-1},
  \label{eq:gaugetrafo}
\end{equation}
where the gauge transformation $\Lambda(t,x)$ is determined by
\begin{equation}
  \dot{\Lambda}=-\lambda \partial_{\nu}B_{\nu}\Lambda,
  \qquad\left.\Lambda\right|_{t=0}=1.
  \label{eq:Lambdatrafo}
\end{equation}
The expectation value of $E$
can thus be computed using the modified
flow equation. Moreover, one is free to set $\lambda=1$, which
turns out to be a particularly convenient choice.

Note that the use of the modified flow equation does not interfere
with the fixing of the gauge of the fundamental field, since
$E$ is unchanged and therefore remains a gauge-invariant
function of the latter.


\subsection{Solution of the modified flow equation}

In perturbation theory the gauge potential is scaled by the
bare coupling,
\begin{equation}
  A_{\mu}\to g_0A_{\mu},
\end{equation}
and the functional integral is then expanded in powers of $g_0$.
The flow $B_{\mu}(t,x)$ thus becomes a function
of the coupling with an asymptotic expansion of the form
\begin{equation}
  B_{\mu}=\sum_{k=1}^{\infty}g_0^kB_{\mu,k},
  \qquad
  \left.B_{\mu,k}\right|_{t=0}=\delta_{k1}A_{\mu}.
  \label{eq:Bseries}
\end{equation}
When this series is inserted in eq.~\eqref{eq:modflow},
and if one sets $\lambda=1$,
a tower of equations
\begin{equation}
  \dot{B}_{\mu,k}-\partial_{\nu}\partial_{\nu}B_{\mu,k}=
  R_{\mu,k},
  \quad k=1,2,\ldots,
  \label{eq:tower}
\end{equation}
is obtained, where the expressions on the right are given by
\begin{align}
  R_{\mu,1}&=0,
  \\
  R_{\mu,2}&=2[B_{\nu,1},\partial_{\nu}B_{\mu,1}]
            -[B_{\nu,1},\partial_{\mu}B_{\nu,1}],
  \\
  R_{\mu,3}&=2[B_{\nu,2},\partial_{\nu}B_{\mu,1}]
            +2[B_{\nu,1},\partial_{\nu}B_{\mu,2}]
            \nonumber\\
  &\hphantom{{}=}
            -[B_{\nu,2},\partial_{\mu}B_{\nu,1}]
            -[B_{\nu,1},\partial_{\mu}B_{\nu,2}]
            +[B_{\nu,1},[B_{\nu,1},B_{\mu,1}]],
\end{align}
and so on. In particular, in $D$ dimensions
the leading-order equation implies
\begin{align}
  B_{\mu,1}(t,x)&=\int\rmd^Dy\,K_t(x-y)A_{\mu}(y),
  \label{eq:B1ker}\\
  K_t(z)&=\int{\rmd^Dp\over(2\pi)^D}\,\rme^{ipz}\rme^{-tp^2}=
  {\rme^{-z^2/4t}\over(4\pi t)^{D/2}},
  \label{eq:Kt}
\end{align}
which shows explicitly that the flow is a smoothing operation.
More precisely, the gauge potential is averaged over a
spherical range in space whose
mean-square radius in four dimensions is
equal to $\sqrt{8t}$.

The higher-order equations~\eqref{eq:tower} can be solved one after another
by noting that
\begin{equation}
   B_{\mu,k}(t,x)=\int_0^t\rmd s\int\rmd^Dy\,K_{t-s}(x-y)R_{\mu,k}(s,y).
   \label{eq:Bksol}
\end{equation}
Recalling eqs.~\eqref{eq:B1ker},\,\eqref{eq:Kt}, it is clear that this
formula generates tree-like expressions, where the fundamental field
$A_{\mu}$ is attached to the endpoints of the trees.

\subsection{Expansion of $\langle E\rangle$}

When the series~\eqref{eq:Bseries} is inserted in
\begin{equation}
\begin{split}
  \langle E\rangle={}&
           \frac{1}{2}\langle\partial_{\mu}B^a_{\nu}\partial_{\mu}B^a_{\nu}-
           \partial_{\mu}B^a_{\nu}\partial_{\nu}B^a_{\mu}\rangle\\
  &+f^{abc}\langle\partial_{\mu}B^a_{\nu}B^b_{\mu}B^c_{\nu}\rangle
           +\frac{1}{4}f^{abe}f^{cde}\langle
           B^a_{\mu}B^b_{\nu}B^c_{\mu}B^d_{\nu}\rangle,
\end{split}
\label{eq:Eexpansion}
\end{equation}
a sequence of terms of increasing order in $g_0$ is obtained.
The lowest-order term is
\begin{equation}
  {\cal E}_0=\frac{1}{2}g_0^2
             \langle\partial_{\mu}B^a_{\nu,1}\partial_{\mu}B^a_{\nu,1}-
             \partial_{\mu}B^a_{\nu,1}\partial_{\nu}B^a_{\mu,1}\rangle
  \label{eq:E0def}
\end{equation}
and the terms at the next order are
\begin{align}
  {\cal E}_1&=g_0^3
    f^{abc}\langle\partial_{\mu}B^a_{\nu,1}B^b_{\mu,1}B^c_{\nu,1}\rangle,
  \\
  {\cal E}_2&=
    g_0^3
    \langle\partial_{\mu}B^a_{\nu,2}\partial_{\mu}B^a_{\nu,1}-
           \partial_{\mu}B^a_{\nu,2}\partial_{\nu}B^a_{\mu,1}\rangle.
  \label{eq:E2def}
\end{align}
Each of these terms is a power series in the gauge coupling,
which may be worked out by expressing the
coefficients $B_{\mu,k}(t,x)$ through the fundamental
field $A_{\mu}(x)$ and by expanding
the correlation functions of the latter using
the standard Feynman rules.

In practice it is advantageous to pass to momentum space by
inserting the Fourier representations
\begin{align}
  B^a_{\mu,1}(t,x)&=\int_p\rme^{ipx}\rme^{-tp^2}\tilde{A}^a_{\mu}(p),
  \label{eq:B1fourier}\\
  B^a_{\mu,2}(t,x)&=if^{abc}\int_0^t\rmd s\int_{q,r}
  \rme^{i(q+r)x}\rme^{-s(q^2+r^2)-(t-s)(q+r)^2}
  \nonumber\\
  &\hphantom{{}=}\times
  \bigl\{\delta_{\mu\lambda}r_{\sigma}-\delta_{\mu\sigma}q_{\lambda}
  +\frac{1}{2}\delta_{\sigma\lambda}(q-r)_{\mu}\bigr\}
  \tilde{A}_{\sigma}^b(q)\tilde{A}_{\lambda}^c(r),
  \label{eq:B2fourier}
\end{align}
and the corresponding
expressions for the higher-order fields $B_{\mu,k}(t,x)$.
The shorthand
\begin{equation}
  \int_{p}=\int{\rmd^D p\over(2\pi)^D}
  \label{eq:intp}
\end{equation}
has been introduced here in order to simplify the notation.

Note that the flow-time integral in eq.~\eqref{eq:B2fourier}
is similar to a Feynman parameter integral.
In particular,
if Feynman parameters are used for
the diagrams contributing to the gluon
correlation functions,
one ends up with gaussian momentum integrals
that can be easily evaluated in any dimension $D$.
The integrals over the parameters then look very much
the same as the ones usually encountered,
except for the fact that the flow-time parameters
are integrated up to $t$ rather than infinity.


\subsection{Computation of ${\cal E}_0$}

In the case of the lowest-order term~\eqref{eq:E0def},
the steps sketched in the previous subsection
lead to the formula
\begin{equation}
  {\cal E}_0=\frac{1}{2}g_0^2(N^2-1)
  \int_p\rme^{-2tp^2}
  (p^2\delta_{\mu\nu}-p_{\mu}p_{\nu})
  D(p)_{\mu\nu},
  \label{eq:E0int}
\end{equation}
where $D(p)_{\mu\nu}$ denotes the unrenormalized full gluon propagator.
Setting $D=4-2\eps$ and choosing the Feynman gauge,
the propagator assumes the form
\begin{align}
  D(p)_{\mu\nu}&={1\over(p^2)^2}\left\{
  (p^2\delta_{\mu\nu}-p_{\mu}p_{\nu})(1-\omega(p))^{-1}+p_{\mu}p_{\nu}
  \right\},
  \\
  \omega(p)&=\sum_{k=1}^{\infty}g_0^{2k}(p^2)^{-k\eps}\omega_k,
\end{align}
from which one infers that
\begin{equation}
  {\cal E}_0=\frac{1}{2}g_0^2{N^2-1\over(8\pi t)^{D/2}}(D-1)
  \left\{1+g_0^2(2t)^{\eps}{\Gamma(2-2\eps)\over\Gamma(2-\eps)}\omega_1
  +\ldots\right\}.
  \label{eq:E0res}
\end{equation}
To this order, the computation is then easily completed by
quoting the known result
\begin{equation}
  \omega_1={1\over 16\pi^2}(4\pi\rme^{-\euler})^{\eps}
  \left\{N\left(\frac{5}{3\eps}+\frac{31}{9}\right)
  -\Nf\left(\frac{2}{3\eps}+\frac{10}{9}\right)
  +\rmO(\eps)\right\}
  \label{eq:omega1}
\end{equation}
for the gluon self-energy, $\euler=0.577\ldots$ being
Euler's constant.


\subsection{Computation of $\langle E\rangle$ to order $g_0^4$}

At the next-to-leading order, the expectation value of $E$
receives contributions
from ${\cal E}_0$ and the terms
${\cal E}_1,{\cal E}_2,\ldots$
up to order $g_0^4$
generated by the expansion of
the flow equation
(cf.~subsect.~2.3). Some of the latter involve
the $3$-point gluon vertex, but in all cases only the
leading-order gluon correlation functions are required.

In the case of the term ${\cal E}_2$, for example,
the purely algebraic part of the calculation leads to
the integral
\begin{equation}
   {\cal E}_2=N(N^2-1)g_0^4\int_0^t\rmd s\int_{q,r}
   {\rme^{-(2t-s)p^2-s(q^2+r^2)}\over p^2q^2r^2}
   \nonumber\\
   \quad\times\bigl\{
   (D-1)p^2(p^2+q^2+r^2)+2(D-2)(q^2r^2-(qr)^2)\bigr\}
   +\rmO(g_0^6)
   \label{eq:E2int}
\end{equation}
where $p=q+r$. This integral appears to be quite complicated, but
the polynomial in the numerator of the integrand allows the
expression to be simplified and eventually to be evaluated
analytically (see Appendix~\ref{app:B} for further details).
The result of the computation,
\begin{equation}
   {\cal E}_2=N(N^2-1){g_0^4\over(4\pi)^D(2t)^{D-2}}
   \bigl\{\frac{9}{2\eps}-\frac{3}{2}+45\ln2-\frac{45}{2}\ln3+
   \rmO(\eps)\bigr\}+\rmO(g_0^6),
   \label{eq:E2res}
\end{equation}
shows that these contributions are not free of ultra-violet
singularities.

The computation of the other terms
follows the same pattern and does not
present any additional difficulties.
Collecting all contributions, the result
\begin{align}
  \langle E\rangle&=\frac{1}{2}g_0^2{N^2-1\over(8\pi t)^{D/2}}(D-1)
  \bigl\{1+c_1g_0^2+\rmO(g_0^4)
  \bigr\},
  \label{eq:Epert}\\
  c_1&={1\over16\pi^2}(4\pi)^{\eps}(8t)^{\eps}\bigl\{
  N\left(\frac{11}{3\eps}+\frac{52}{9}-3\ln3\right)
  -\Nf\left(\frac{2}{3\eps}+\frac{4}{9}-\frac{4}{3}\ln2\right)
  +\rmO(\eps)\bigr\},
\end{align}
is then obtained.


\subsection{Renormalization}

The bare coupling $g_0$ is related to
the renormalized coupling $g$
in the $\MSbar$ scheme
and the associated normalization mass $\mu$
according to~\cite{BardeenEtAl}
\begin{align}
  g_0^2&=g^2\mu^{2\eps}(4\pi\rme^{-\euler})^{-\eps}
  \Bigl\{1-\frac{1}{\eps}b_0g^2+\rmO(g^4)\Bigr\},
  \label{eq:g0rel}\\
  b_0&={1\over16\pi^2}\left\{\frac{11}{3}N-\frac{2}{3}\Nf\right\}.
\end{align}
Now if eq.~\eqref{eq:g0rel} is written as an expansion in the renormalized coupling,
the terms proportional to $1/\eps$ cancel and one obtains
\begin{align}
  \langle E\rangle&={3(N^2-1)g^2\over128\pi^2t^2}
  \bigl\{1+\bar{c}_1g^2+\rmO(g^4)
  \bigr\},
  \label{eq:Eren}\\
  \bar{c}_1&=
  {1\over16\pi^2}\bigl\{
  N\left(\frac{11}{3}L+\frac{52}{9}-3\ln3\right)
  -\Nf\left(\frac{2}{3}L+\frac{4}{9}-\frac{4}{3}\ln2\right)\bigr\},
\end{align}
at $\eps=0$, where $L=\ln(8\mu^2t)+\euler$.
To this order in the gauge coupling,
$\langle E\rangle$ thus turns out to be a
quantity that does not need to be renormalized
and which therefore encodes some physical property of
the theory.

In terms of the running coupling $\alpha(q)$
at scale $q=(8t)^{-1/2}$,
the expansion~\eqref{eq:Eren} assumes the form
\begin{align}
  \langle E\rangle&={3(N^2-1)\over32\pi t^2}\alpha(q)
  \bigl\{1+k_1\alpha(q)+\rmO(\alpha^2)
  \bigr\},
  \label{eq:Erun}\\
  k_1&=
  {1\over4\pi}\bigl\{
  N\left(\frac{11}{3}\euler+\frac{52}{9}-3\ln3\right)
  -\Nf\left(\frac{2}{3}\euler+\frac{4}{9}-\frac{4}{3}\ln2\right)\bigr\}.
\end{align}
In particular, for $N=3$ one obtains
\begin{equation}
  \langle E\rangle={3\over4\pi t^2}\alpha(q)
  \bigl\{1+k_1\alpha(q)+\rmO(\alpha^2)
  \bigr\},
  \qquad
  k_1=1.0978+0.0075\times\Nf.
  \label{eq:E3}
\end{equation}
The next-to-leading order correction is reasonably small in
this case and corresponds
to a change in the momentum $q$ by a factor $2$ or so if $\Nf\leq3$.

